\section{Related Work}

Edge caching has been widely investigated across wireless networks, mobile systems, and IoT-enabled edge computing. Existing work can be broadly categorized into (i) optimization-based cooperative caching, (ii) mobility-aware and social-aware caching, (iii) reinforcement learning–based caching, and (iv) knowledge and graph-based caching approaches. Despite significant progress, current approaches still exhibit limitations in explicitly modeling dynamic, mobility-derived topology as a first-class learning object.

\subsection{Optimization-Based Cooperative Caching}

Several works formulate cooperative caching as an optimization problem under storage and bandwidth constraints. Blasco and Gunduz~\cite{Blasco2014} proposed a learning-based proactive caching framework that models content popularity dynamics using online learning and solves cache placement via convex optimization. Similarly, Maddah-Ali and Niesen~\cite{MaddahAli2014} introduced coded caching to improve network efficiency through coordinated placement and delivery strategies. These approaches demonstrate that coordination among nodes can significantly improve system performance. However, they typically assume static or centrally known network topology and do not adaptively learn the mobility-induced relationships between nodes. Xu et al. proposed a cooperative caching framework for UAV-assisted vehicular networks that combines LSTM-attention-based popularity prediction with coalition-game-based cache optimization, making it most closely related to optimization-based cooperative caching \cite{bai2025collaborative}. Ma et al. proposed a delay-oriented caching strategy for LEO satellite-assisted edge caching networks by modeling a content request popularity with a Zipf distribution and average transmission delay using stochastic geometry, and by optimizing satellite caching probabilities through a particle swarm optimization based minimum-delay strategy, which improves delay performance through probabilistic cache allocation, but it does not explicitly represent mobility-driven relationships among edge nodes or use topology-aware graph learning for cooperative caching decisions \cite{ma2025delay}.

\subsection{Mobility-Aware and Social-Aware Caching}

User mobility has been recognized as an important factor in caching efficiency. Wang et al.~\cite{Wang2017Mobility} explored mobility prediction to improve caching decisions in mobile edge networks. Mardham et al.~\cite{Mardham2018} proposed opportunistic distributed caching in delay-tolerant networks using mobility patterns. Cabaniss and Madria~\cite{Cabaniss2014} propose a socially-aware caching strategy for DTNs that leverages contact frequency and social relationships to guide content placement. Their approach improves delivery probability by exploiting community structures but relies on predefined social metrics rather than explicit topology modeling.
While these works acknowledge the correlations, cooperation is often driven by predefined mobility predictions or social heuristics rather than by learning the evolving topological structure among edge nodes.

\subsection{Reinforcement Learning-Based Edge Caching}

Reinforcement learning (RL) has emerged as a promising technique for dynamic caching under non-stationary demand. DQN-style methods use replay buffers and target networks to stabilize learning over discrete actions \cite{mnih2015dqn,vanhasselt2016double}. Wang et al.~\cite{Min2020MARL} proposed a multi-agent deep reinforcement learning framework MARL with bandwidth-efficient inter-agent communication to improve cooperative edge caching performance under dynamic demand. Although it introduces cooperative behavior, the underlying topology is typically implicit. Agents exchange limited information or coordinate through reward shaping, but the structural graph of inter-node relationships is not explicitly modeled or learned as a dynamic entity. Jiang et al. proposed cooperative content caching in vehicular edge computing networks as a Markov decision process and proposed a two-stage deep reinforcement learning framework for joint caching decisions across vehicles and roadside units to reduce long-term delay and improve cache hit rate under dynamic vehicular mobility, heterogeneous node constraints, and a large caching policy space \cite{jiang2026cooperative}. Dhara et al. proposed a caching framework for post-disaster Space-Air-Ground Integrated Networks (SAGIN) that uses a Contextual Multi-Armed Bandit and a Federated MAB model to support real-time content retrieval across satellites, UAVs, base stations, and vehicles. The framework addresses dynamic caching under heterogeneous network conditions, limited storage, mobility, and changing request patterns by using contextual features such as content popularity, freshness, and cache-size ratio into caching decisions, while the federated extension improves adaptation by sharing learned reward information among distributed nodes \cite{dhara2026co}. Cao et al. proposed a mobility-aware cooperative caching framework for vehicular edge computing based on federated distillation and deep reinforcement learning which predicts popular content using mobility information and privacy-aware distributed learning, and then uses reinforcement learning to decide how that content should be cached across neighboring roadside units \cite{cao2025}. Zhao et al. \cite{zhao2026collaborative} proposed a collaborative caching and offloading method for vehicular edge computing that predicts the stability of links between vehicles, analyzes the popularity of historical offloading tasks for caching, and then uses deep reinforcement learning to make partial offloading decisions.

\subsection{Knowledge and Graph Based Caching Approaches}

More recently, structured representations such as knowledge graphs have been introduced to enhance caching decisions. Wang et al.~\cite{wang2024knowledge} proposed a knowledge graph based reinforcement learning framework for cooperative caching in mobile edge computing systems. Their approach captures semantic relationships among contents and network states. 

Despite incorporating structured representations, most existing work focuses on content level relationships rather than explicitly modeling time varying inter node mobility graphs. The integration of graph learning methods to directly encode mobility-induced topology in cooperative caching remains largely unexplored.
